\scanpage{16}
Soit $p:X\to Y$ un morphisme submersif de schémas formels (lisses), \unsure{supposons} $X$ muni d'une $\mathcal A_*$-structure. On dit que $Y$ admet une $\mathcal A_*$-structure quotient de celle de $X$ (moyennant $p$) s'il admet une \inserted{$\mathcal A_*$} structure telle que $p$ devienne un $\mathcal A_*$-morphisme. Celle-ci est encore unique: en effet, il existe (sans condition sur $Y$) un \deleted{\unsure{schéma}} sous-schéma $\mathcal A_*$-formel $Y'$ de $X$, tel que le morphisme formel induit $Y'\to Y$ soit un isomorphisme. \deleted{Le} \inserted{Considérons le} composé
\diagram{\begin{tikzcd}[column sep=small,row sep=small]&Y\arrow[d,"\sim"]\\X\arrow[ur,dash,"p"]\arrow[r]&Y'\end{tikzcd}}
\deleted{\unsure{de $Y'$}}, on voit \unsure{dès qu'à cs|dès grâce à ce} que la $\mathcal A_*$-structure quotient existe ssi $\unsure{X'|X}\to Y'$ est un $\mathcal A_*$-morphisme, et la $\mathcal A_*$-structure sur $Y$ est celle déduite de celle de $Y'$.

\deleted{\unsure{Un cor}} Soit $X\xrightarrow{p}Y$ un morphisme formel submersif. Alors $R=X\mathbin{\hat\times}_Y X\subset X\mathbin{\hat\times}X$ est une \inserted{sous-schéma formel lisse, et c'est} relation d'équivalence [stable par symétrie, contient diagonale, transitif $\ldots$]. \deleted{\unsure{Inversement}} \inserted{Plus précisément} $X$ étant donné, il y a corr. biunivoque entre ces quotients formels $Y$, et les $R$ rel. d'équivalence formelles. Si $X$ est muni d'une $\mathcal A_*$-structure, alors si $Y$ admet une $\mathcal A_*$-structure quotient, on voit que $R$ est un sous-schéma $\mathcal A_*$-formel de $X\mathbin{\hat\times}X$ i.e. il a une $\mathcal A_*$-structure induite.

f) (Condition intégrale) La réciproque est vraie, i.e. si $R$ est une relation d'équival. formelle dans $X$, qui soit un \unsure{$\mathcal A_*$-sous-structure|$\mathcal A_*$-sous-schéma} de $X\mathbin{\hat\times}X$, alors
\mspage{15}
